% Source: physical PDF pp. 136--145 (printed pp. 113--122), from the
% Section 27.1 heading through the paragraph immediately before Section 27.2.
% Inventory: Eqs. (27.1.1)--(27.1.28), ten unnumbered display groups, one
% linked citation, and no footnotes, figures, tables, or typographic dividers.
% Uncertainties: none.

\subsection{Gauge-Invariant Actions for Chiral Superfields}

Consider a set of Abelian or non-Abelian gauge transformations that leave
the supersymmetry generator \(Q\) invariant. (For simple supersymmetry
there is just one Majorana spinor supersymmetry generator, which can only
furnish a trivial representation of any semi-simple gauge group.) Each
component field in a supermultiplet must transform in the same way under
such gauge transformations. In particular, for a left-chiral superfield we
have
\begin{equation}
  \begin{aligned}
  \phi_n(x)&\longrightarrow
    \sum_m\left[
      \exp\left(\ii\sum_A t_A\Lambda^A(x)\right)
    \right]_{nm}\phi_m(x)\,,
  \\
  \psi_{nL}(x)&\longrightarrow
    \sum_m\left[
      \exp\left(\ii\sum_A t_A\Lambda^A(x)\right)
    \right]_{nm}\psi_{mL}(x)\,,
  \\
  \mathcal{F}_n(x)&\longrightarrow
    \sum_m\left[
      \exp\left(\ii\sum_A t_A\Lambda^A(x)\right)
    \right]_{nm}\mathcal{F}_m(x)\,,
  \end{aligned}
  \tag{27.1.1}\label{eq:27.1.1}
\end{equation}
where \(t_A\) are Hermitian matrices representing the generators of the
gauge algebra, and \(\Lambda^A(x)\) are real functions of \(x^\mu\) that
parameterize a finite gauge transformation. (We are using the same notation
for gauge transformations as in Section 15.1, except that in order to avoid
confusion with Dirac indices we label the gauge generators and gauge
transformation parameters with the letters \(A,B\), etc.\ instead of
\(\alpha,\beta\), etc.)

The left-chiral superfield \eqref{eq:26.3.11} involves derivatives of some
of the component fields, so its transformation is more complicated than
Eq.~\eqref{eq:27.1.1}. However, Eq.~\eqref{eq:26.3.21} shows that the
superfield does not involve derivatives if expressed in terms of
\(\theta_L\) and the variable \(x_+\) defined by Eq.~\eqref{eq:26.3.23}.
It therefore has the transformation property
\begin{equation}
  \Phi_n(x,\theta)\longrightarrow
  \sum_m\left[
    \exp\left(\ii\sum_A t_A\Lambda^A(x_+)\right)
  \right]_{nm}\Phi_m(x,\theta)\,.
  \tag{27.1.2}\label{eq:27.1.2}
\end{equation}
If a term in the action depends only on left-chiral superfields, and not
their derivatives or complex conjugates, like the term
\(\int \dd^4x\,[f(\Phi)]_{\mathcal{F}}\) in Eq.~\eqref{eq:26.3.30}, then it
(and its complex conjugate) will be invariant under the local
transformation \eqref{eq:27.1.2} if it is invariant under global
transformations with \(\Lambda^A(x)\) independent of \(x^\mu\). The need
for introducing gauge fields in renormalizable theories of chiral
superfields arises only in the \(D\)-terms, which involve both \(\Phi_n\)
and \(\Phi_n^*\). Because the matrices \(t_A\) are Hermitian, the Hermitian
adjoint of Eq.~\eqref{eq:27.1.2} is
\begin{equation}
  \Phi_n^\dagger(x,\theta)\longrightarrow
  \sum_m\Phi_m^\dagger(x,\theta)
  \left[
    \exp\left(-\ii\sum_A t_A\Lambda^A(x_+)^*\right)
  \right]_{mn}\,.
  \tag{27.1.3}\label{eq:27.1.3}
\end{equation}
If it were not for the difference between
\(\Lambda^A(x_+)^*=\Lambda^A(x_-)\) and \(\Lambda^A(x_+)\), this would
just say that \(\Phi^\dagger\) transforms according to the representation
of the gauge group that is contragredient to the representation furnished
by \(\Phi\), and any function of \(\Phi\) and \(\Phi^\dagger\) that is
invariant under global gauge transformations would also be invariant under
local gauge transformations. Because \(x_+\) and \(x_-\) are different, we
must introduce a gauge connection matrix \(\Gamma_{nm}(x,\theta)\), with
the transformation property
\begin{equation}
  \begin{split}
  \Gamma(x,\theta)\longrightarrow{}&
  \exp\left(+\ii\sum_A t_A\Lambda^A(x_+)^*\right)
  \Gamma(x,\theta)
  \exp\left(-\ii\sum_A t_A\Lambda^A(x_+)\right).
  \end{split}
  \tag{27.1.4}\label{eq:27.1.4}
\end{equation}
Then by multiplying \(\Phi^\dagger\) on the right with \(\Gamma\) we
obtain a superfield that transforms as
\begin{equation}
  \begin{split}
  \left[\Phi^\dagger(x,\theta)\Gamma(x,\theta)\right]_n
  \longrightarrow
  \sum_m\left[\Phi^\dagger(x,\theta)\Gamma(x,\theta)\right]_m
  \left[
    \exp\left(-\ii\sum_A t_A\Lambda^A(x_+)\right)
  \right]_{mn},
  \end{split}
  \tag{27.1.5}\label{eq:27.1.5}
\end{equation}
so that any globally gauge-invariant function constructed from \(\Phi\)
and \(\Phi^\dagger\Gamma\) (and not their derivatives or complex
conjugates) will also be locally gauge-invariant. One obvious example is
the gauge-invariant version \((\Phi^\dagger\Gamma\Phi)_D\) of the
\(D\)-term in the Lagrangian constructed in Section 26.4.

Any \(\Gamma(x,\theta)\) that transforms as in Eq.~\eqref{eq:27.1.4} will
allow us to construct gauge-invariant Lagrangians of chiral superfields.
The choice is not unique; if \(\Gamma\) transforms as in
Eq.~\eqref{eq:27.1.4}, and we multiply on the right with any left-chiral
superfield \(\Upsilon_L\) with the transformation rule
\[
  \Upsilon_L(x,\theta)\longrightarrow
  \exp\left(\ii\sum_A t_A\Lambda^A(x_+)\right)
  \Upsilon_L(x,\theta)
  \exp\left(-\ii\sum_A t_A\Lambda^A(x_+)\right),
\]
then we obtain a new gauge connection that also satisfies
Eq.~\eqref{eq:27.1.4}. One simplification is to take
\(\Gamma(x,\theta)\) Hermitian:
\begin{equation}
  \Gamma^\dagger(x,\theta)=\Gamma(x,\theta)\,.
  \tag{27.1.6}\label{eq:27.1.6}
\end{equation}
This is always possible if there is any \(\Gamma(x,\theta)\) that satisfies
Eq.~\eqref{eq:27.1.4}, for then by taking the Hermitian adjoint of
Eq.~\eqref{eq:27.1.4} we easily see that \(\Gamma^\dagger(x,\theta)\)
transforms in the same way as \(\Gamma(x,\theta)\), so if
\(\Gamma(x,\theta)\) is not Hermitian then we can replace it with its
Hermitian part \((\Gamma+\Gamma^\dagger)/2\) (or, if that vanishes, by its
anti-Hermitian part \((\Gamma-\Gamma^\dagger)/2\ii\).) Another simplification
of great physical importance is to express \(\Gamma(x,\theta)\) in terms
of fields whose gauge transformation properties do not depend on the
specific representation \(t_A\) of the gauge algebra under which the
chiral superfield \(\Phi(x,\theta)\) transforms, so that these fields can
be used to form a suitable matrix \(\Gamma(x,\theta)\) for chiral
superfields that transform according to any representation of the gauge
group. For this purpose, it is useful to recall the Baker--Hausdorff
formula, which states that, for arbitrary matrices \(a\) and \(b\),
\begin{equation}
  e^a e^b
  =\exp\left(
    a+b+\frac12[a,b]
    +\frac1{12}[a,[a,b]]
    +\frac1{12}[b,[b,a]]
    +\cdots
  \right),
  \tag{27.1.7}\label{eq:27.1.7}
\end{equation}
where `\(\cdots\)' denotes higher-order terms that can be written as
multiple commutators of \(a\)s and \(b\)s, like the second- and third-order
terms that are shown explicitly. It follows from this that for any
representation of a Lie algebra, we have
\begin{equation}
  \exp\left(\sum_A a^A t_A\right)
  \exp\left(\sum_A b^A t_A\right)
  =\exp\left(\sum_A f^A(a,b)t_A\right),
  \tag{27.1.8}\label{eq:27.1.8}
\end{equation}
where
\begin{equation}
  \begin{aligned}
  f^A(a,b)={}&a^A+b^A
  +\frac{\ii}{2}\sum_{BC}C^A{}_{BC}a^B b^C
  -\frac1{12}\sum_{BCDE}C^A{}_{BC}C^C{}_{DE}a^B a^D b^E
  \\
  &-\frac1{12}\sum_{BCDE}C^A{}_{BC}C^C{}_{DE}b^B b^D a^E
  +\cdots,
  \end{aligned}
  \tag{27.1.9}\label{eq:27.1.9}
\end{equation}
which depends on the Lie algebra through its structure constants
\(C^A{}_{BC}\), defined as usual by
\[
  [t_B,t_C]=\ii\sum_A C^A{}_{BC}t_A\,,
\]
but does not depend on the particular representation furnished by the
\(t_A\). We will therefore take \(\Gamma(x,\theta)\) in the form
\begin{equation}
  \Gamma(x,\theta)
  =\exp\left(-2\sum_A t_A V^A(x,\theta)\right),
  \tag{27.1.10}\label{eq:27.1.10}
\end{equation}
where \(V^A(x,\theta)\) are a set of real superfields (so that \(\Gamma\)
is Hermitian), not depending on the representation of the gauge algebra
furnished by the \(t_A\).

We can achieve an important further simplification by noting an additional
symmetry of supersymmetric gauge theories. If some function of \(\Phi\)
and \(\Phi^\dagger\Gamma\) is invariant under global gauge transformations,
then it will automatically be invariant not only under the local gauge
transformations \eqref{eq:27.1.2}--\eqref{eq:27.1.4}, but also under the
larger group of \textit{extended gauge transformations}
\begin{equation}
  \Phi_{nL}(x,\theta)\longrightarrow
  \sum_m\left[
    \exp\left(\ii\sum_A t_A\Omega^A(x,\theta)\right)
  \right]_{nm}\Phi_{mL}(x,\theta)
  \tag{27.1.11}\label{eq:27.1.11}
\end{equation}
and
\begin{equation}
  \begin{split}
  \Gamma(x,\theta)\longrightarrow{}&
  \exp\left(-\ii\sum_A t_A\Omega^A(x,\theta)\right)
  \Gamma(x,\theta)
  \exp\left(+\ii\sum_A t_A\Omega^A(x,\theta)^*\right),
  \end{split}
  \tag{27.1.12}\label{eq:27.1.12}
\end{equation}
where \(\Omega^A(x,\theta)\) is an arbitrary left-chiral superfield---that
is, an arbitrary function of \(\theta_L\) as well as \(x_+\). Under this
transformation,
\begin{equation}
  V^A(x,\theta)\longrightarrow
  V^A(x,\theta)
  +\frac{\ii}{2}\left[
    \Omega^A(x,\theta)-\Omega^A(x,\theta)^*
  \right]+\cdots,
  \tag{27.1.13}\label{eq:27.1.13}
\end{equation}
where `\(\cdots\)' denotes terms arising from the commutators in
Eq.~\eqref{eq:27.1.7}, which are of first or higher order in gauge coupling
constants. As a general left-chiral superfield, \(\Omega\) may be written
in the form \eqref{eq:26.3.11}
\begin{equation}
  \begin{aligned}
  \Omega^A(x,\theta)={}&
  W^A(x)
  -\sqrt{2}\left(
    \bar\theta\frac{1+\gamma_5}{2}w^A(x)
  \right)
  +\mathcal{W}^A(x)\left(
    \bar\theta\frac{1+\gamma_5}{2}\theta
  \right)
  \\
  &+\frac12(\bar\theta\gamma_5\gamma_\mu\theta)
    \partial^\mu W^A(x)
  -\frac1{\sqrt{2}}(\bar\theta\gamma_5\theta)
    \left(
      \bar\theta\sl{\partial}\frac{1+\gamma_5}{2}w^A(x)
    \right)
  \\
  &-\frac18(\bar\theta\gamma_5\theta)^2\Box W^A(x),
  \end{aligned}
  \tag{27.1.14}\label{eq:27.1.14}
\end{equation}
in which \(W^A(x)\) and \(\mathcal{W}^A(x)\) are arbitrary complex
functions of \(x^\mu\), and we introduced Majorana spinors \(w^A(x)\)
defined so that the left-handed spinor components of the superfields are
\(\tfrac12(1+\gamma_5)w^A(x)\). Using the complex conjugation properties
\eqref{eq:26.A.21} of Majorana bilinears, the complex conjugate of
Eq.~\eqref{eq:27.1.14} gives
\begin{equation}
  \begin{aligned}
  \Omega^A(x,\theta)^*={}&
  W^{A*}(x)
  -\sqrt{2}\left(
    \bar\theta\frac{1-\gamma_5}{2}w^A(x)
  \right)
  +\mathcal{W}^{A*}(x)\left(
    \bar\theta\frac{1-\gamma_5}{2}\theta
  \right)
  \\
  &-\frac12(\bar\theta\gamma_5\gamma_\mu\theta)
    \partial^\mu W^{A*}(x)
  -\frac1{\sqrt{2}}(\bar\theta\gamma_5\theta)
    \left(
      \bar\theta\sl{\partial}\frac{1-\gamma_5}{2}w^A(x)
    \right)
  \\
  &-\frac18(\bar\theta\gamma_5\theta)^2\Box W^{A*}(x).
  \end{aligned}
  \tag{27.1.15}\label{eq:27.1.15}
\end{equation}
We write the real superfields \(V^A(x,\theta)\) in terms of component
fields as in Eq.~\eqref{eq:26.2.10}:
\begin{equation}
  \begin{aligned}
  V^A(x,\theta)={}&C^A(x)
  -\ii\bigl(\bar\theta\gamma_5\omega^A(x)\bigr)
  -\frac{\ii}{2}(\bar\theta\gamma_5\theta)M^A(x)
  -\frac12(\bar\theta\theta)N^A(x)
  \\
  &+\frac{\ii}{2}(\bar\theta\gamma_5\gamma^\mu\theta)V^A_\mu(x)
  -\ii(\bar\theta\gamma_5\theta)
    \left(
      \bar\theta\left[\lambda^A(x)
      +\frac12\sl{\partial}\omega^A(x)\right]
    \right)
  \\
  &-\frac14(\bar\theta\gamma_5\theta)^2
    \left(D^A(x)+\frac12\Box C^A(x)\right),
  \end{aligned}
  \tag{27.1.16}\label{eq:27.1.16}
\end{equation}
where \(C^A(x)\), \(M^A(x)\), \(N^A(x)\), and \(V^A_\mu(x)\) are all
real, and \(\omega^A(x)\) and \(\lambda^A(x)\) are Majorana spinors.
Using Eqs.~\eqref{eq:27.1.14}--\eqref{eq:27.1.16} in
Eq.~\eqref{eq:27.1.13}, we find that the component fields of the gauge
superfield undergo the extended gauge transformation
\begin{equation}
  \begin{aligned}
  C^A(x)&\longrightarrow C^A(x)-\operatorname{Im}W^A(x)+\cdots,
  \\
  \omega^A(x)&\longrightarrow
    \omega^A(x)+\frac1{\sqrt{2}}w^A(x)+\cdots,
  \\
  V^A_\mu(x)&\longrightarrow
    V^A_\mu(x)+\partial_\mu\operatorname{Re}W^A(x)+\cdots,
  \\
  M^A(x)&\longrightarrow
    M^A(x)-\operatorname{Re}\mathcal{W}^A(x)+\cdots,
  \\
  N^A(x)&\longrightarrow
    N^A(x)+\operatorname{Im}\mathcal{W}^A(x)+\cdots,
  \\
  \lambda^A(x)&\longrightarrow\lambda^A(x)+\cdots,
  \\
  D^A(x)&\longrightarrow D^A(x)+\cdots,
  \end{aligned}
  \tag{27.1.17}\label{eq:27.1.17}
\end{equation}
where again `\(\cdots\)' denotes terms that arise from the structure
constants in Eq.~\eqref{eq:27.1.9} and that are therefore proportional
to one or more factors of gauge coupling constants. We can use such an
extended gauge transformation to put the gauge superfields into a
convenient form, known as \textit{Wess--Zumino gauge}
~\hyperref[ch27-ref-1]{[1]}, in which
\begin{equation}
  C^A(x)=\omega^A(x)=M^A(x)=N^A(x)=0\,,
  \tag{27.1.18}\label{eq:27.1.18}
\end{equation}
so that
\begin{equation}
  \begin{aligned}
  V^A(x,\theta)={}&
  \frac{\ii}{2}(\bar\theta\gamma_5\gamma^\mu\theta)V^A_\mu(x)
  -\ii(\bar\theta\gamma_5\theta)\bigl(\bar\theta\lambda^A(x)\bigr)
  -\frac14(\bar\theta\gamma_5\theta)^2D^A(x).
  \end{aligned}
  \tag{27.1.19}\label{eq:27.1.19}
\end{equation}
To accomplish this to zeroth order in the coupling constants, it is only
necessary to set \(\operatorname{Im}W^A(x)=C^A(x)\),
\(w^A(x)=-\sqrt{2}\omega^A(x)\), and
\(\mathcal{W}^A(x)=M^A(x)-\ii N^A(x)\). For Abelian gauge theories, in which
the structure constants vanish, this ends our task. For non-Abelian gauge
theories it is necessary to add terms to \(\operatorname{Im}W^A(x)\),
\(w^A(x)\), and \(\mathcal{W}^A(x)\) of first order in gauge coupling
constants to cancel the terms arising from commutators of zeroth-order
terms, then add terms to \(\operatorname{Im}W^A(x)\), \(w^A(x)\), and
\(\mathcal{W}^A(x)\) of second order in gauge coupling constants to cancel
the terms arising from commutators of first-order terms with zeroth-order
terms, and so on. It is not easy to calculate the series of terms in
\(\operatorname{Im}W^A(x)\), \(w^A(x)\), and \(\mathcal{W}^A(x)\) needed
to satisfy the gauge conditions \eqref{eq:27.1.18} to all orders in gauge
couplings, but there is no need---the important thing is that it is possible.

Inspection of the transformation rules
\eqref{eq:26.2.11}--\eqref{eq:26.2.14} shows that the Wess--Zumino gauge
condition \eqref{eq:27.1.18} is not invariant under supersymmetry
transformations unless \(V^A_\mu=\lambda^A=0\), and the condition
\(\lambda^A=0\) is not supersymmetric unless also \(D^A=0\), in which
case the whole superfield vanishes. Once we adopt Wess--Zumino gauge, the
action is no longer invariant under either general extended gauge
transformations or under supersymmetry, but it is invariant under
supersymmetry transformations, which take us out of Wess--Zumino gauge,
followed by suitable extended gauge transformations that take us back to
Wess--Zumino gauge. (We will go into this explicitly in Section 27.8.) As
we shall now see, it is also invariant under the ordinary gauge
transformations \eqref{eq:27.1.2}--\eqref{eq:27.1.4}, which preserve
Wess--Zumino gauge.

With the gauge superfield satisfying the Wess--Zumino gauge condition
\eqref{eq:27.1.18}, it becomes relatively easy to calculate its behavior
under ordinary infinitesimal gauge transformations. In this case,
\(\Omega^A(x_+)\) are left-chiral superfields of the form
\eqref{eq:26.3.11}, but with no \(\psi_L\)- or
\(\mathcal{F}\)-components, and with \(\phi\)-components given by
\textit{real} infinitesimal functions \(\Lambda^A(x)\):
\begin{equation}
  \Omega^A(x_+)=\Lambda^A(x)
  +\frac12(\bar\theta\gamma_5\gamma_\mu\theta)
    \partial^\mu\Lambda^A(x)
  -\frac18(\bar\theta\gamma_5\theta)^2\Box\Lambda^A(x).
  \tag{27.1.20}\label{eq:27.1.20}
\end{equation}
To calculate the product of exponentials in the transformation rule
\eqref{eq:27.1.4}, we use a version of the Baker--Hausdorff formula:
\begin{equation}
  \exp(a)\exp(X)\exp(b)
  =\exp\left[
    X+L_X\mathbin{\cdot}(b-a)
    +(L_X\coth L_X)\mathbin{\cdot}(b+a)+\cdots
  \right],
  \tag{27.1.21}\label{eq:27.1.21}
\end{equation}
where \(a\), \(b\), and \(X\) are arbitrary matrices, \(L_X\) is the
operator
\begin{equation}
  L_X\mathbin{\cdot} f=\frac12[X,f]\,,
  \tag{27.1.22}\label{eq:27.1.22}
\end{equation}
and `\(\cdots\)' here denotes terms of second and higher order in \(a\)
and/or \(b\). In our case we have
\[
  \begin{aligned}
  b+a
  &=2\sum_A t_A\operatorname{Im}\Lambda^A(x_+)
   =-\ii(\bar\theta\gamma_5\gamma_\mu\theta)
     \sum_A t_A\partial^\mu\Lambda^A(x),
  \\
  b-a
  &=-2\ii\sum_A t_A\operatorname{Re}\Lambda^A(x_+)
   =-2\ii\sum_A t_A\left[
     \Lambda^A(x)-\frac18(\bar\theta\gamma_5\theta)^2
     \Box\Lambda^A(x)
   \right],
  \\
  X
  &=-2\sum_A t_AV^A(x,\theta)
  \\
  &=-2\sum_A t_A\left[
     \frac{\ii}{2}(\bar\theta\gamma_5\gamma^\mu\theta)V^A_\mu(x)
     -\ii(\bar\theta\gamma_5\theta)(\bar\theta\lambda^A(x))
     -\frac14(\bar\theta\gamma_5\theta)^2D^A(x)
   \right].
  \end{aligned}
\]
Now, every term in \(X\) involves at least one factor of \(\theta_L\) and
at least one factor of \(\theta_R\), while \(a+b\) has just one factor of
\(\theta_L\) and one factor of \(\theta_R\), so we can drop any terms in
\(L_X\coth L_X\) of second or higher order in \(L_X\). Since
\(L_X\coth L_X\) is an \textit{even} function of \(L_X\), this means that
we can replace it with its term of zeroth order in \(L_X\), which is just
unity. Also, we may drop the term in \(b-a\) proportional to
\((\bar\theta\gamma_5\theta)^2\), since when acted on by \(L_X\) it would
yield at least three factors of either \(\theta_L\) or \(\theta_R\). Thus
the argument of the exponential on the right-hand side of
Eq.~\eqref{eq:27.1.21} may be replaced with
\[
  \begin{aligned}
  X+\frac12[X,b-a]+b+a
  &=-2\sum_A t_A\Bigl[V^A(x,\theta)
  \\
  &\hspace{22mm}
    +\sum_{BC}C^A{}_{BC}V^B(x,\theta)\Lambda^C(x)
  \\
  &\hspace{22mm}
    +\frac{\ii}{2}(\bar\theta\gamma_5\gamma_\mu\theta)
      \partial^\mu\Lambda^A(x)\Bigr].
  \end{aligned}
\]
Thus for infinitesimal gauge transformations, the transformation rule
\eqref{eq:27.1.4} yields
\begin{equation}
  \begin{split}
  V^A(x,\theta)\longrightarrow
  V^A(x,\theta)
  +\sum_{BC}C^A{}_{BC}V^B(x,\theta)\Lambda^C(x)
  +\frac{\ii}{2}(\bar\theta\gamma_5\gamma_\mu\theta)
    \partial^\mu\Lambda^A(x).
  \end{split}
  \tag{27.1.23}\label{eq:27.1.23}
\end{equation}
It is important to note that under ordinary gauge transformations, a gauge
superfield in Wess--Zumino gauge remains in Wess--Zumino gauge. In terms of
the component fields in Eq.~\eqref{eq:27.1.19},
Eq.~\eqref{eq:27.1.23} reads
\begin{equation}
  V^A_\mu(x)\longrightarrow
  \sum_{BC}C^A{}_{BC}V^B_\mu(x)\Lambda^C(x)
  +\partial_\mu\Lambda^A(x),
  \tag{27.1.24}\label{eq:27.1.24}
\end{equation}
\begin{equation}
  \lambda^A(x)\longrightarrow
  \sum_{BC}C^A{}_{BC}\lambda^B(x)\Lambda^C(x),
  \tag{27.1.25}\label{eq:27.1.25}
\end{equation}
\begin{equation}
  D^A(x)\longrightarrow
  \sum_{BC}C^A{}_{BC}D^B(x)\Lambda^C(x).
  \tag{27.1.26}\label{eq:27.1.26}
\end{equation}
We recognize Eq.~\eqref{eq:27.1.24} as the usual Yang--Mills gauge
transformation rule (15.1.9) of a gauge field, while
Eqs.~\eqref{eq:27.1.25} and \eqref{eq:27.1.26} tell us that the fields
\(\lambda^A(x)\) and \(D^A(x)\) transform as `matter' fields that belong
to the adjoint representation of the gauge group. The Majorana spinors
\(\lambda^A\) are known as \textit{gaugino} fields, while the real scalars
\(D^A\) will turn out to be another set of auxiliary fields.

Next we must evaluate the matrix \(\Gamma\) that is needed in constructing
gauge-invariant functions of chiral superfields. Because all terms with more
than four factors of \(\theta\) vanish, the expansion of the exponential is
very simple in Wess--Zumino gauge:
\[
  \begin{aligned}
  \Gamma(x,\theta)
  &=\exp\left(-2\sum_A t_AV^A(x,\theta)\right)
  \\
  &=\mathbf{1}
  -\ii(\bar\theta\gamma_5\gamma^\mu\theta)
    \sum_A t_AV^A_\mu(x)
  \\
  &\quad-\frac12(\bar\theta\gamma_5\gamma^\mu\theta)
    (\bar\theta\gamma_5\gamma^\nu\theta)
    \sum_{AB}t_At_BV^A_\mu(x)V^B_\nu(x)
  \\
  &\quad+2\ii(\bar\theta\gamma_5\theta)
    \sum_A t_A(\bar\theta\lambda^A(x))
  +\frac12(\bar\theta\gamma_5\theta)^2
    \sum_A t_AD^A(x).
  \end{aligned}
\]
We can construct a gauge-invariant density by multiplying this on the right
with a column vector of left-chiral superfields of the form
\eqref{eq:26.3.11}:
\[
  \begin{aligned}
  \Phi_n(x,\theta)={}&
  \phi_n(x)-\sqrt{2}(\bar\theta\psi_{nL}(x))
  +\mathcal{F}_n(x)\left(
    \bar\theta\frac{1+\gamma_5}{2}\theta
  \right)
  \\
  &+\frac12(\bar\theta\gamma_5\gamma_\mu\theta)
    \partial^\mu\phi_n(x)
  -\frac1{\sqrt{2}}(\bar\theta\gamma_5\theta)
    (\bar\theta\sl{\partial}\psi_{nL}(x))
  -\frac18(\bar\theta\gamma_5\theta)^2\Box\phi_n(x),
  \end{aligned}
\]
and multiplying on the left with the column
\[
  \begin{aligned}
  \Phi_n(x,\theta)^*={}&
  \phi_n^*(x)-\sqrt{2}(\overline{\psi_{nL}(x)}\theta)
  +\mathcal{F}_n^*(x)\left(
    \bar\theta\frac{1-\gamma_5}{2}\theta
  \right)
  \\
  &-\frac12(\bar\theta\gamma_5\gamma_\mu\theta)
    \partial^\mu\phi_n^*(x)
  -\frac1{\sqrt{2}}(\bar\theta\gamma_5\theta)
    \partial_\mu\bigl(\overline{\psi_{nL}(x)}\gamma^\mu\theta\bigr)
  \\
  &-\frac18(\bar\theta\gamma_5\theta)^2\Box\phi_n^*(x).
  \end{aligned}
\]
The term in this product of fourth order in \(\theta\) is
\[
  \begin{aligned}
  \left[\Phi^\dagger\Gamma\Phi\right]_{\theta^4}
  ={}&-\frac18(\bar\theta\gamma_5\theta)^2
  \left\{
    [\phi^\dagger\Box\phi]+[(\Box\phi^\dagger)\phi]
  \right\}
  \\
  &+(\bar\theta\gamma_5\theta)
  \left\{
    [(\overline{\psi_L}\theta)(\bar\theta\gamma^\mu
      \partial_\mu\psi_L)]
    +[((\partial_\mu\overline{\psi_L})\gamma^\mu\theta)
      (\bar\theta\psi_L)]
  \right\}
  \\
  &+\frac14\bigl(\bar\theta(1-\gamma_5)\theta\bigr)
    \bigl(\bar\theta(1+\gamma_5)\theta\bigr)
    [\mathcal{F}^\dagger\mathcal{F}]
  \\
  &-\frac14(\bar\theta\gamma_5\gamma^\mu\theta)
    (\bar\theta\gamma_5\gamma^\nu\theta)
    [\partial_\mu\phi^\dagger\partial_\nu\phi]
  \\
  &-\frac{\ii}{2}(\bar\theta\gamma_5\gamma^\mu\theta)
    (\bar\theta\gamma_5\gamma^\nu\theta)
    \sum_A V^A_\mu
    \left\{
      [\phi^\dagger t_A\partial_\nu\phi]
      -[(\partial_\nu\phi^\dagger)t_A\phi]
    \right\}
  \\
  &-\frac12(\bar\theta\gamma_5\gamma^\mu\theta)
    (\bar\theta\gamma_5\gamma^\nu\theta)
    \sum_{AB}V^A_\mu V^B_\nu[\phi^\dagger t_At_B\phi]
  \\
  &-2\ii(\bar\theta\gamma_5\gamma_\mu\theta)
    \sum_A V^{A\mu}
    [(\overline{\psi_L}\theta)t_A(\bar\theta\psi_L)]
  \\
  &-2\ii\sqrt{2}(\bar\theta\gamma_5\theta)
    \sum_A[(\overline{\psi_L}\theta)t_A
      (\bar\theta\lambda^A)\phi]
  \\
  &-2\ii\sqrt{2}(\bar\theta\gamma_5\theta)
    \sum_A[\phi^\dagger(\overline{\lambda^A}\theta)t_A
      (\bar\theta\psi_L)]
  \\
  &+\frac12(\bar\theta\gamma_5\theta)^2
    \sum_A D_A[\phi^\dagger t_A\phi],
  \end{aligned}
\]
in which we use square brackets to indicate scalar products in the flavor
indices \(n,m\), and continue to use round brackets to indicate scalar
products in Dirac indices. Just as in Section 26.4, we may use the
identities \eqref{eq:26.A.17}--\eqref{eq:26.A.19} to put all \(\theta\)
dependence in the form of an over-all factor
\((\bar\theta\gamma_5\theta)^2\):
\[
  \begin{aligned}
  \left[\Phi^\dagger\Gamma\Phi\right]_{\theta^4}
  =(\bar\theta\gamma_5\theta)^2\biggl\{&
  -\frac18[\phi^\dagger\Box\phi]
  -\frac18[(\Box\phi^\dagger)\phi]
  \\
  &+\frac14[(\overline{\psi_L}\gamma^\mu\partial_\mu\psi_L)]
  -\frac14[((\partial_\mu\overline{\psi_L})\gamma^\mu\psi_L)]
  \\
  &-\frac12[\mathcal{F}^\dagger\mathcal{F}]
  +\frac14[\partial_\mu\phi^\dagger\partial^\mu\phi]
  \\
  &+\frac{\ii}{2}\sum_A V^A_\mu
    [\phi^\dagger t_A\partial^\mu\phi]
  -\frac{\ii}{2}\sum_A V^A_\mu
    [(\partial^\mu\phi^\dagger)t_A\phi]
  \\
  &+\frac12\sum_{AB}V^A_\mu V^{B\mu}
    [\phi^\dagger t_At_B\phi]
  -\frac{\ii}{2}\sum_A V^A_\mu
    [(\overline{\psi_L}\gamma^\mu t_A\psi_L)]
  \\
  &-\frac{\ii}{\sqrt{2}}\sum_A
    [(\overline{\psi_L}t_A\lambda^A)\phi]
  +\frac{\ii}{\sqrt{2}}\sum_A
    [\phi^\dagger(\overline{\lambda^A}t_A\psi_L)]
  \\
  &+\frac12\sum_A D_A[\phi^\dagger t_A\phi]
  \biggr\}.
  \end{aligned}
\]
The \(D\)-term is the coefficient of
\(-\tfrac14(\bar\theta\gamma_5\theta)^2\) minus \(\tfrac12\Box\) acting
on the \(\theta\)-independent term, which for
\([\Phi^\dagger\Gamma\Phi]\) is \([\phi^\dagger\phi]\), so
\[
  \begin{aligned}
  [\Phi^\dagger\Gamma\Phi]_D={}&
  -2[\partial_\mu\phi^\dagger\partial^\mu\phi]
  -[(\overline{\psi_L}\gamma^\mu\partial_\mu\psi_L)]
  +[((\partial_\mu\overline{\psi_L})\gamma^\mu\psi_L)]
  +2[\mathcal{F}^\dagger\mathcal{F}]
  \\
  &-2\ii\sum_A V^A_\mu[\phi^\dagger t_A\partial^\mu\phi]
  +2\ii\sum_A V^A_\mu[(\partial^\mu\phi^\dagger)t_A\phi]
  \\
  &-2\sum_{AB}V^A_\mu V^{B\mu}
    [\phi^\dagger t_At_B\phi]
  +2\ii\sum_A V^A_\mu
    [(\overline{\psi_L}\gamma^\mu t_A\psi_L)]
  \\
  &+2\ii\sqrt{2}\sum_A[(\overline{\psi_L}t_A\lambda^A)\phi]
  -2\ii\sqrt{2}\sum_A[\phi^\dagger(\overline{\lambda^A}t_A\psi_L)]
  \\
  &-2\sum_A D_A[\phi^\dagger t_A\phi].
  \end{aligned}
\]
To see that this is gauge-invariant, we note that it may be written as
\begin{equation}
  \begin{aligned}
  \frac12[\Phi^\dagger\Gamma\Phi]_D={}&
  -[(D_\mu\phi)^\dagger D^\mu\phi]
  -\frac12[(\overline{\psi_L}\gamma^\mu D_\mu\psi_L)]
  +\frac12[(\overline{D_\mu\psi_L}\gamma^\mu\psi_L)]
  +[\mathcal{F}^\dagger\mathcal{F}]
  \\
  &+\ii\sqrt{2}\sum_A[(\overline{\psi_L}t_A\lambda^A)\phi]
  -\ii\sqrt{2}\sum_A[\phi^\dagger(\overline{\lambda^A}t_A\psi_L)]
  -\sum_A D_A[\phi^\dagger t_A\phi],
  \end{aligned}
  \tag{27.1.27}\label{eq:27.1.27}
\end{equation}
where \(D_\mu\) is the gauge-invariant derivative (15.1.10):
\begin{equation}
  D_\mu\psi_L\equiv
  \partial_\mu\psi_L-\ii\sum_A t_AV^A_\mu\psi_L,
  \qquad
  D_\mu\phi\equiv
  \partial_\mu\phi-\ii\sum_A t_AV^A_\mu\phi.
  \tag{27.1.28}\label{eq:27.1.28}
\end{equation}
Eq.~\eqref{eq:27.1.27} is thus a suitable gauge-invariant kinematic
Lagrangian for the scalar and spinor components of a left-chiral
superfield, now supplemented with Yukawa couplings of gaugino fields to
the scalar and spinor components of chiral superfields as well as terms
involving auxiliary fields \(\mathcal{F}_n\) and \(D_A\).
